% generated by roam.py handbook export -- не редактировать, править заметки
% 20261004152534-виды_градиентного_спуска.org (body 42d7cb5f0033)

\begin{concept}[{Виды градиентного спуска}]\label[concept]{note:BEEE81AD-13FC-4846-A160-C72E31760BEE}
  Рассмотрим общую схему метода спуска из точки $x_k$:
  \[
  x_{k+1} = x_k + \alpha_kd_k
  \]
  Где: $d_k$~--- направление спуска,  $\alpha_k$~--- шаг спуска.

  Как выбрать шаг спуска:

  \begin{enumerate}
    \item Фиксированный шаг
    \item \hyperref[note:8EDDABAC-D2BE-4EC9-8297-DBAC52F70D28]{Точная одномерная оптимизация} $\alpha_k = \arg\min_{\alpha > 0} f(x_k + \alpha d_k)$
    \item \hyperref[note:A818191F-B1A6-492C-B38D-5369D169161C]{Неточная одномерная оптимизация}: не обязательно на каждом шаге искать идеальное решение. Пример: \hyperref[note:5330B23D-95DD-4D5D-B4A3-63EBADB1DB7D]{условия Армихо-Вульфа}.
  \end{enumerate}
\end{concept}

\begin{remark}
  При $d_k = -\nabla f(x_k)$ это градиентный спуск.
\end{remark}
